\documentclass[10pt,a4paper]{amsart}
\usepackage[margin=19mm]{geometry}
\usepackage{amsmath,amssymb,mathtools}
\usepackage[scr=rsfs,cal=euler]{mathalfa}
\usepackage{xfrac}
\usepackage[unicode,hidelinks,bookmarks=false]{hyperref}
\newcommand{\ie}{i.e.\ }
\newcommand{\cf}{cf.\ }
\newcommand{\resp}{respectively\ }
\newcommand{\Subsection}{\subsection}
\providecommand{\nobreakdash}{\nobreak\hbox{-}}
\setlength{\emergencystretch}{2em}
\multlinegap0pt
\usepackage{amssymb}
\usepackage[all]{xy}
\usepackage{tikz}
\usepackage{breqn}
\newtheorem{proposition}{Proposition}
\def\T{{\mathcal T}}
\newcommand{\p}{{\partial}}
\title{Genus zero Whitham hierarchy via Hurwitz--Frobenius manifolds}
\author{Alexey Basalaev}
\date{Source: arXiv:2507.07615v2 (CC BY 4.0)}
\begin{document}
\maketitle

\noindent\emph{Source and licence.} Genus zero Whitham hierarchy via Hurwitz--Frobenius manifolds. Version arXiv:2507.07615v2 (CC BY 4.0), distributed by arXiv under the Creative Commons Attribution 4.0 International licence, http://creativecommons.org/licenses/by/4.0/, as recorded in the arXiv licence metadata for this exact version and shown on the abstract page https://arxiv.org/abs/2507.07615v2. Full licence notice: \texttt{licenses/arxiv-2507.07615-CC-BY-4.0.txt}.\par\smallskip

\noindent\textit{Editorial scope.} Counted unit: the article's proposition stabilization (source0.tex lines 721--727). The article's complete original proof of it is delivered with it (source0.tex lines 728--736, one \texttt{proof} environment of 9 lines). No mathematics is written, repaired or re-derived: the statement and the proof are the article's own lines, in the article's own notation, and no retained line is substituted or re-flowed.  Retained uncounted as the source-local context the proof needs: lines 515--521.  Nothing was omitted from any retained span, so the package carries no omission record.  No retained line contains a citation, so the wrapper carries no bibliography and no bibliography entry.  No retained line contains a figure environment, an image-inclusion command or drawing code, so no image is delivered and the package declares no asset dependency.  Editorial formatting only, in addition: an AMS \texttt{amsart} preamble that restates the article's own declarations and macros used by the retained lines, namely the environments \texttt{proposition} on separate counters, together with the standard packages those lines need.  The retained environments print in this document as Proposition 1 in order of appearance, whereas the article prints the same items with the numbering of its own full document.  The complete native source of the article as distributed in the arXiv e-print for this exact version is bundled unchanged as \texttt{sources/181359/source0.tex}.
\begin{proposition}\label{prop: low ramification products}
  Let $X,Y,Z \in \T_M$ be such that $(X \cdot \lambda)(Y \cdot \lambda) = (Z \cdot \lambda)$
%   \[
%      (X \cdot \lambda)(Y \cdot \lambda) = (Z \cdot \lambda)
%   \]
    assumed as the equality of rational functions. Then $X \circ Y = Z$.
\end{proposition}

\begin{proposition}\label{prop: v-product stabilization}
    Let $L \ge K$, $i \neq j$ and $a < k_i$, $b < k_j$. Then in the basis $ \frac{\p}{\p v}$
    \[
        {}^{(K)}c_{i,a;j,b}^{k,l} = {}^{(L)}c_{i,a;j,b}^{k,l} \quad \forall k,l
    \]
    assumed at the equality of functions in $v_{j,\alpha}$.
\end{proposition}

\begin{proof}
    The statement is equivalent to the claim that
    \[
        \frac{1}{(z - q_i)^{a+1}} \circ \frac{1}{(z - q_j)^{b+1}}
    \]
    has the same expression via $\p \lambda / \p v_{k,l}$ in both Hurwitz--Frobenius manifolds. Because $i \neq j$ the rational function $\frac{1}{(z - q_i)^{a+1}} \cdot \frac{1}{(z - q_j)^{b+1}}$ belongs to both Hurwitz--Frobenius manifolds assumed and can therefore be expressed via $\p \lambda / \p v_{k,l}$.    
    
    The statement follows now by Proposition~\ref{prop: low ramification products} because in the $v$--coordinates the Hurwitz space $M_K$ is embedded into $M_L$ just by setting some of the coordinates to zero.
\end{proof}

\end{document}
